\section{Ramsey Theory}
Just for sake of completeness, we define
\[
	R(k,l) = \min\{ n : K_k \not \subset G \subset K_n \Rightarrow
	\alpha(G) \geq l \}.
\]
In coloring terms, it's the smallest $n$ such that any red and blue
coloring of $E(K_n)$ contains a red $K_k$ or a blue $K_l$.

\subsection{Bounds on \mbox{R(3,k)}}
We have two very simple bounds that are easy to prove,
\[
	\bigg(\frac{k}{\log(k)}\bigg)^{3/2} \leq R(3,k) \leq \binom{k+1}{2}.
\]
The upperbound follows easily from the Erdös-Szekeres proof.
\begin{theorem}
	$R(3,k) \leq \binom{k+1}{2}$
\end{theorem}
\begin{proof}

	Let $G$ on $n$-vtx be a triangle-free with no independet set of size $k$,
	we know every neighborhood is an independent set so $\Delta(G) < k$.
	Now we greedily get one vtx at a time removing their neighbors.
	Pick $v_1$ and $N(v_1)$ with $|N(v_1)| \leq k-1$, then pick $v_2$,
	notice that the neighborhood of $v_2$ (after removing $v_1$ and its neighbors)
	has size at most $k-2$, if not $\{v_1\} \cup N(v_2)$ would be a
	independent set of size $k$. We proceed like this for $v_1, v_2, \dots v_t$,
	where each step we discard at most $k-(i-1)$ vtxs. Because we can't do this
	process $k$ times, for some $i < k$ we must fail (meaning there are
	no vertex left),
	so $n = |V(G)| \leq \sum_{j=0}^{k-2} k - j \leq \binom{k+1}{2} - 1$.
\end{proof}

The lowerbound is a beautiful application of the probabilistic
method. The main idea is to choose $G(n,p)$
as to have few triangles (linearly small in $n$) and relatively low
independence number.

\begin{theorem}
	\[
		R(3,k) \geq \bigg(\frac{k}{\log(k)}\bigg)^{3/2}
	\]
\end{theorem}

We need a \textbf{very important} lemma first.
\begin{lemma}
	With high probability, the independence number of $G(n,p) \leq
	2\log(n)/p$. That is
	\[
		\Pr(\alpha(G(n,p)) \geq 2\log(n)/p + 1) = 1 - o_{(n,p(n))}(1)
	\]
	where $p$ can be taken polynomially small in $n$.
\end{lemma}
\begin{proof}
	\begin{align*}
		\Pr(\alpha(G(n,p)) \geq k) &\leq \Ex[\#\text{independent k sets}
		\in G(n,p)]\\
		&= \binom{n}{k}(1-p)^{\binom{k}{2}} \leq
		\bigg(\frac{en}{k}\bigg)^k\cdot e^{-pk(k-1)/2}\\
		&= \bigg(\frac{en}{e^{p(k-1)/2}k}\bigg)^k
	\end{align*}
	Taking $k \geq 2\log(n)/p + 1$ we get that
	\[
		\Pr(\alpha(G(n,p)) \geq k) \leq
		\bigg(\frac{ep}{2\log(n)}\bigg)^{2\log(n)/p} \to 0
	\]
	very fast even! (faster than any polynomial).
\end{proof}

The strategy will be to remove triangles from $G(n,p)$ by removing
one vertex from each. This cannot increase
the independence number, so we hope the resulting graph still has
many vertices to yield a Ramsey construction.

So let's see what this alteration may give us.
\begin{proof}
	(Of the Theorem) We choose $p$ small such that the number of
	triangles is no larger than a small linear in $n$ (say $n/2$).
	\[
		\Ex[\# K_3 \subset G(n,p)] = \binom{n}{3}p^3 = (1 + o(1))\frac{(pn)^3}{6}
	\]
	We want
	\[
		\Pr( \# K_3 \in G(n,p) \geq n/2) < (1 + o(1))\frac{p^3n^2}{3} < 1/2
	\]
	so taking $p = \eps n^{-2/3}$ for some small $\eps$ we get the
	inequality. But what does this $p$ gives in independence number?
	We have w.h.p (say bigger than one half)
	\[
		\alpha(G(n,p)) \leq 2\eps^{-1}\log(n)n^{2/3}.
	\]
	So $\Pr(K_3(G(n,p)) < n/2) > 1/2$ and $\Pr(\alpha(G(n,p) <
	C\log(n)n^{2/3})) > 1/2$, meaning there exists
	$G$ satisfying both statements. Remove one vtx from each triangle
	in $G$, we are left with a triangle-free graph $G'$ with at least $n/2$ vtxs
	and $\alpha(G) < C\log(n)n^{2/3}$.
	Setting $k = C\log(n)n^{2/3}$, we get that
	\[
		V(G') = n/2 \geq c\bigg(\frac{k}{\log(k)}\bigg)^{3/2}
	\]
\end{proof}

Now for a better (and astounishingly beautiful) bound.
\begin{theorem}
	\label{trm:krivilevich_r3k}
	\[
		R(3,k) \geq \bigg(\frac{k}{\log(k)}\bigg)^2.
	\]
\end{theorem}
The idea of this bound is removing edges. Maybe we can kill the
triangles by removing few edges and not
increase the independence number too much. For this we will need a
beautiful lemma.

\begin{lemma}
	\label{lemma:erdos_tetali}
	(Erdös-Tetali) Let $\{A_1, A_2, \dots A_m\}$ be a family of events
	in some probability space $(\Omega, P)$. Let
	$\mathcal{E}_t$ be the event that at least $t$ different
	independent $A_i$ happen. We have that
	\[
		\Pr(\mathcal{E}_t) \leq \frac{1}{t!}\bigg(\sum_{i=1}^{m} A_i\bigg)^t.
	\]
\end{lemma}
\begin{proof}
	We get this one from just opening the terms
	\[
		\Pr(\mathcal{E}_t) = \Pr\bigg(\bigcup_{\substack{i_1 < \dots <
		i_t\\ A_{i_j} \not\sim A_{i_k}}} \bigcap_{j = 1}^{t} A_{i_j}\bigg)
		\leq \sum_{\substack{i_1 < \dots < i_t\\ A_{i_j} \not\sim
		A_{i_k}}} \prod_{j = 1}^t \Pr(A_{i_j}) \leq \sum_{\substack{i_1
		< \dots < i_t}} \prod_{j = 1}^t \Pr(A_{i_j}) \leq
		\frac{1}{t!}\bigg(\sum_{i=1}^{m} A_i\bigg)^t.
	\]
\end{proof}

The second ingredient is noting that if you're a bit bigger (say a
constant times) than the independent size, you already have many
edges (at least half of what is expected).

\begin{lemma}
	\label{lemma:concentration_edges}
	Let $k \geq C\log(n)/p$, then
	\[
		\Pr\bigg(\exists \text{ k-set } S, e(G(n,p)[S]) <
		p\binom{k}{2}/2\bigg) = o(1)
	\]
	for some $C > 0$.
\end{lemma}
\begin{proof}
	Let $\mathcal{E}$ be the event that there exists a k-set $S$ with
	$e(G(n,p)[S]) < p\binom{k}{2}/2$.
	We will do a union bound and use Hoeffding (or Chernoff) to bound
	its probability.

	Let $S$ be any $k$-set. Note that (letting $B(n,p)$ be the binomial
	with parameters $n$ and $p$) we have
	\[e(G(n,p)[S]) \equiv B\bigg(\binom{k}{2},p\bigg). \]
	So by Hoeffding and its bounds on the binomial
	[\ref{lemma:expec_binomial_bounds}] (need to add
	proofs for these inequalities),
	\begin{align*}
		\Pr\bigg(e(G(n,p)) < \frac{p}{2}\binom{k}{2}\bigg) &=
		\Pr\bigg(\bigg|B\bigg(\binom{k}{2}, p\bigg) - p\binom{k}{2}\bigg|
		\geq \frac{p}{2}\binom{k}{2}\bigg)
		\leq e^{-c p\binom{k}{2}}
	\end{align*}
	for some $c > 0$.

	Now, doing the union bound over $k$-sets we find that (taking $C > 3c^{-1}$)
	\begin{align*}
		\Pr(\mathcal{E}) \leq \binom{n}{k}e^{-c p\binom{k}{2}} \leq
		\bigg(\frac{en}{e^{cp(k-1)/2}k}\bigg)^k \leq
		\bigg(\frac{e}{k}\bigg)^k \to 0.
	\end{align*}
\end{proof}

The idea to prove theorem \ref{trm:krivilevich_r3k} is to remove a
maximal collection of edge-disjoint triangles from $G(n,p)$, (notice
that this makes it triangle free).
We use lemma \ref{lemma:erdos_tetali} to show that this procedure
won't remove too many edges from any set $S$ with size $C\log(n)/p$.

\begin{proof}
	(Of Theorem \ref{trm:krivilevich_r3k}) Let $S$ be a set of size $k
	= C\log(n)/p$, for some $p$ we will determine later.
	Let $\{T_1, \dots T_m\}$ be the set of all triangles that intersect
	with at least a pair in $S$. Note that $m \leq n\binom{|S|}{2}$.
	Let $\mathcal{E}_S$ be the event that a edge disjoint set of these
	triangles with size bigger or equal to $\frac{p}{6}\binom{k}{2}$
	are in $G(n,p)$.
	By [\ref{lemma:erdos_tetali}], we can bound $\Pr(\mathcal{E}_S)$,

	\[
		\Pr(\mathcal{E}_S) \leq
		\frac{1}{[\frac{p}{6}\binom{k}{2}]!}\bigg(\sum_{i = 1}^{m}
		\Pr(T_i \in G(n,p))\bigg)^t \leq
		\bigg(\frac{24e}{pk^2}\bigg)^{pk^2/24}(p^3nk^2)^{pk^2/12}
	\]
	Replacing $k$ in the equation we get (and putting $C > 24e$)
	\[
		\Pr(\mathcal{E}_S) \leq
		\bigg(\frac{24ep}{C\log^2(n)}\bigg)^{\log(n)^2/(24p)}
		(pn\log(n)^2)^{\log(n)^2/(12p)}
		\leq (p^3n^2\log(n)^2)^{\log(n)^2/(24p)}
	\]
	So, if we take $p = \eps n^{-1/2}$ (might try to optimize more
	later), we already find that
	\[
		\Pr(\mathcal{E}_S) \leq
		\bigg(\frac{\log(n)}{n^{1/2}}\bigg)^{n^{1/2}\log^2(n)} \ll
		\bigg(\frac{C\log(n)}{n^{1/2}}\bigg)^{C\log(n)n^{1/2}} \leq
		\binom{n}{k}^{-1}
	\]
	So $\mathcal{E}_S$ does not happen for any $S$ with high
	probability. As w.h.p $S$ of size $k$ have more than
	$\frac{p}{2}\binom{k}{2}$ edges, and
	no maximal disjoint triangle family intersect too much with $S$. We
	can remove a maximal edge disjoint triangle family and not make any
	set of size $k$ independent.
	This yields a triangle free graph on $n$ vtxs with independence
	number less than $C\log(n)n^{1/2}$, so
	\[
		R(3,k) \geq c\bigg(\frac{k}{\log k}\bigg)^2.
	\]
\end{proof}

Now let's find an upperbound of the right order.
\begin{theorem}
	\label{trm:r3k_aks}
	(Ajtai, Komlós, Szemerédi)
	\[
		R(3,k) \leq C\frac{k^2}{\log k}
	\]
\end{theorem}

The theorem will be immediately implied by the following lemma.
\begin{lemma}
	\label{lemma:independence_k3_free}
	Let $G$ be $K_3$-free with $\Delta(G) \leq d$, then
	\[
		\alpha(G) \geq \frac{cn}{d}\log(d).
	\]
	for some universal constant $c > 0$.
\end{lemma}

\begin{proof}
	(Of the theorem assuming Lemma \ref{lemma:independence_k3_free})
	Notice that $G$ is $K_3$ free and
	its independence number is less than $k$, then, as every
	neighborhood is a independent set, we have $\Delta(G) < k$.
	Applying the lemma we get that
	\[
		k > \alpha(G) \geq \frac{cn}{k}\log(k)
	\]
	which implies
	\[
		n \leq C\frac{k^2}{\log(k)}.
	\]
\end{proof}

The proof of the lemma is exceptionally fun. The main idea is to take
a uniform random independent set in $G$,
and show that in expectation, its size is bigger than
$\frac{cn}{d}\log(d)$ - which is very unnintuitive.
\begin{proof}
	(Of Lemma \ref{lemma:independence_k3_free}) Let $I \in
	\mathcal{I}(G)$ be a uniformly chosen random independent set of $G$.
	For each $v \in V(G)$, define the random variable
	\[
		X_v = d \cdot \id[v \in I] + \sum_{u \in N(v)} \id[u \in I].
	\]
	Note that
	\[
		\Ex \sum_{v \in V(G)} X_v = d\cdot \Ex |I| + \sum_{v \in V(G)}
		\sum_{u \in N(v)} \Pr(u \in I) =
		d\cdot \Ex |I| + \sum_{u \in V(G)} d(u) \Pr(u \in I) \leq 2d\cdot \Ex |I|.
	\]
	So we have that $\Ex |I| \geq \frac{1}{2d} \Ex [\sum_{v \in V(G)}
	X_v]$. To prove the lemma,
	it would suffice to show then that $\Ex X_v \geq \log(d)$ for all
	$v$ (as summing it would yield the result).

	So let's lowerbound $\Ex X_v$. The idea is to condition it on
	something and find a bound that works \textbf{FOR EVERY} condition
	we take (which seems insane).

	Fix $v$, let's divide the world into $U = \{v\} \cup N(v)$ and $W =
	V(G)\setminus U$ (both depend on $v$). Notice
	that the value of $X_v$ only changes for elements in $U$ that
	appear in $I$. So fix an independent set $S$ of $W$, and let's
	calculate $\Ex[X_v | I \cap W = S]$.

	From the fact that $I$ has to be an independent set, we know it
	cannot have any neighbors
	of elements in $S$. So let $Y = N(v) \cap N(S)$ and let
	$X = N(v) \setminus Y$ (note $v \not \in N(S)$). As $I$ is
	uniformly chosen, we now know that $I \cap U$ (conditioned on $I \cap W = S$),
	is a uniformly chosen independent set of $\{v\} \cup X$. Let $|X| =
	t$, remembering that $X$ is an independent set, we find
	\begin{align*}
		\Ex[X_v | I \cap W = S] &= \Ex[X_v | v \in I, I \cap W =
		S]\cdot\Pr(v\in I | I \cap W = S)\\
		&+ \Ex[X_v | v \not \in I, I \cap W = S]\cdot\Pr(v \not \in I | I
		\cap W = S)\\
		&= \frac{1}{2^t + 1}d + \frac{2^t}{2^t + 1} \frac{t}{2}
	\end{align*}
	where we used the fact that $I \cap U$ is an uniform independent
	set of the star $G[\{v\} \cup X]$. Now (here comes the crazy part),
	notice that independent of what $t = |X|$ is, we have
	\[
		\frac{1}{2^t + 1}d + \frac{2^t}{2^t + 1} \frac{t}{2} \geq \frac{\log(d)}{8}.
	\]
	just take $t < \log(d)/2$ and $t \geq \log(d)/2$. But this means
	\[
		\Ex X_v = \Ex[\Ex[X_v| I \cap W(v)]] \geq \frac{\log(d)}{8}
	\]
	yielding the result.
\end{proof}

\subsection{Bounds on R(4,k)}

By doing vertex deletion and Erdös-Szekeres, we may easily find (we
need $p^6n^4 < n$, so $p = \eps n^{-1/2}$)
\[
	c\bigg(\frac{k}{\log(k)}\bigg)^2\leq R(4,k) \leq k^3.
\]
Doing the Erdös-Tetali method [\ref{lemma:erdos_tetali}] (we need
$p^6n^4 < pn^2$, so $p = \eps n^{-2/5}$)
we get also
\[
	c\bigg(\frac{k}{\log k}\bigg)^{5/2} \leq R(4,k).
\]
But we can do better on the upperbound by doing a similar strategy as
in the $R(3,k)$ case. The next lemma Rob stated,
but I don't think we will use it.

\begin{lemma}
	Let $G$ on $n$ vtx be $K_4$-free and $\Delta(G) \leq d$, then
	\[
		\alpha(G) \geq \frac{cn}{d}\frac{\log (d)}{\log\log(d)}.
	\]
\end{lemma}

\begin{proof}
	The proof of this is similar to lemma
	\ref{lemma:independence_k3_free}, we must however pick a uniformly independent
	set from the now $K_3$-free $X$. (I will try to do the computations
	more carefully later)
\end{proof}

To get an even better bound, we will need some kind of stability
setting in these super sparse graphs.
\begin{lemma}
	\label{lemma:AKSz}
	(AKSz) Let $G$ be a graph on $n$-vtx with $\Delta(G) \leq d$ and
	$k_3(G) \leq (nd^2)/\lambda$,
	then
	\[
		\alpha(G) \geq c\frac{n}{d}\log(\lambda)
	\]
	where $c > 0$ is an absolute constant.
\end{lemma}

With it we can prove the following strong bound,

\begin{theorem}
	\label{trm:aksz_r4k}
	\[
		R(4,k) \leq C\frac{k^3}{\log(k)^2}.
	\]
\end{theorem}
\begin{proof}
	(Assuming [\ref{lemma:AKSz}]) Let $G$ be $n$-vtx, $K_4$-free graph
	with $\alpha(G) < k$,
	we know by Erdös-Szekeres that $\Delta(G) < R(3,k)$. How many
	triangles can $G$ have?
	Pick a vertex $v$ and let $S = N(v)$. We know that for any vertex
	$u \in S$, $N_{G[S]}(u)$ is independent
	(because $G$ is $K_4$-free) and so $d_{G[S]}(u) < k$. This implies that
	\[
		k_3(G) = \frac{1}{3} \sum_{v \in V(G)} e(G[N(v)]) \leq \frac{nk^2}{6}
	\]
	by our bound on $\Delta(G[S])$.

	So applying [\ref{lemma:AKSz}] with $d = R(3,k)$ and $\lambda =
	(k/\log(k))^2$ yields
	\[
		k > \alpha(G) > c\frac{n\log(k)^2}{k^2}
	\]
	which yields $n < Ck^3/\log(k)^2$.
\end{proof}

So now it suffices to prove \ref{lemma:AKSz}.
\begin{proof}
	(Of the AKSz lemma)
	The idea will be to do some vertex deletion and apply
	\ref{lemma:independence_k3_free}, but we don't really
	know how so we do it randomly.

	To easy up the calculation (it only changes the constant in front),
	suppose $\Delta(G) \leq d$ and $k_3(G) \leq nd^2/\lambda^3$.
	Take a $q$-random subset $Q$ of $V(G)$, we can estimate some
	statistics, we find:
	\begin{enumerate}[label=(\roman*)]
		\item $\Ex |Q| = qn$
		\item $\Ex e(G[Q]) = q^2e(G) \leq q^2nd/2$
		\item $\Ex k_3(G[Q]) = q^3k_3(G) \leq (q/\lambda)^3nd^2$
	\end{enumerate}

	By taking $q = \lambda/d$ and using Markov to remove triangles and
	high degree vtxs, we find a set of  size at least $qn/6$ with no triangles
	and $\Delta(G[Q']) < 3q\Delta(G)$. Applying
	\ref{lemma:independence_k3_free} to this set yields the result.

	(Formally)
	We have very good concentration on $|Q|$ (as long as $qn \to
	\infty$), so applying Chernoof and Markov (twice) we
	see that
	\begin{align*}
		\Pr\bigg( \bigg(|Q| > \frac{qn}{2}\bigg) &\land \bigg(e(G[Q]) <
			3q^2e(G)\bigg)
		\land \bigg(k_3(G[Q]) < \frac{3q^3nd^2}{\lambda^3}\bigg) \bigg)\\
		&= 1 -\Pr\bigg( \bigg(|Q| < \frac{qn}{2}\bigg) \lor \bigg(e(G[Q])
			> 3q^2e(G)\bigg)
		\lor \bigg(k_3(G[Q]) > \frac{3q^3nd^2}{\lambda^3}\bigg) \bigg)\\
		&\geq 1 - o(1) - 1/3 - 1/3 > 0
	\end{align*}

	So we find a $G' = G[Q]$ for some good $Q$ satisfying all these
	properties. $G'$ satisfies:
	\begin{enumerate}[label=(\roman*)]
		\item $v(G') \geq \frac{qn}{2}$
		\item $e(G') \leq 3q^2e(G) \leq 3q^2nd/2$
		\item $k_3(G') \leq \frac{3q^3nd^2}{\lambda^3}$
	\end{enumerate}

	By removing at most half of the vertices of $G'$ (using Markov) and
	killing a vertex from each triangle,
	we find triangle free $G''$ with $v(G') \geq \frac{qn}{4} - k_3(G')
	\geq \frac{qn}{8}$ and $\Delta(G'') \leq 6q^2nd/2$.
\end{proof}

This seems like a very interesting strategy for off-diagonal Ramsey.
We did the case $l = 3$,
$l=4$, so a very natural conjecture would be:
\begin{theorem}
	\[
		R(l,k) \leq C\frac{k^{l-1}}{(\log k)^{l-2}}
	\]
\end{theorem}
\begin{proof}
	We prove it by induction. The case $l = 3,4$ were already done.
	Let's follow the proof of [\ref{trm:aksz_r4k}].
	Suppose $G$ is an $n$-vtx, $K_l$-free graph with $\alpha(G) < k$.
	We must have $\Delta(G) < R(l-1,k) \leq d$. How
	many triangles can it have? We know that if $u \in N(v)$, then
	$|N(u)\cap N(v)| < R(l-2,k)$. So, doing the same bound as before,
	we see that $k_3(G) \leq ndR(l-2,k)/3$. So, to use \ref{lemma:AKSz}
	we must have
	\[
		\frac{ndR(l-2,k)}{3} \leq \frac{nd^2}{\lambda}
	\]
	so
	\[
		\lambda \leq \frac{3d}{R(l-2,k)}
	\]
	and (by induction), it suffices that $\lambda$ satisfies
	\[
		\lambda \leq \frac{c d(\log k)^{l-2}}{k^{l-1}}
	\]
	applying \ref{lemma:AKSz} with $d = \frac{k^{l-1}}{(\log k)^{l-2}}
	\geq R(l-1,k) > \Delta(G)$, we have by hypohesis,
	\[
		\alpha(G) \geq \frac{n}{d}\log(\lambda) \geq C_l
		\frac{n\log(k)^{l-2}}{k^{l-2}}
	\]
	and therefore,
	\[
		n \leq C_l\frac{k^{l-1}}{(\log k)^{l-2}}
	\]
\end{proof}

This proof can also be found in Rob's ICM manuscript.
[\url{https://arxiv.org/pdf/2601.05221}].

\subsection{Bounds on R(3,3,k)}
So what are some easy bounds we can get on this? An easy lowerbound
is $R(3,k)$, by doing Erdös-Szekeres,
we can get the multinomial bound as well
\[
	R(3,k) \leq R(3,3,k) \leq 2kR(3,k) \leq C\frac{k^3}{\log k}.
\]
We can do a bit better on the upperbound by counting triangles and
applying [\ref{lemma:AKSz}]!
\begin{theorem}
	\[R(3,3,k) \leq \frac{k^3}{(\log k)^2}\]
\end{theorem}
\begin{proof}
	Suppose there is a coloring $\chi$ on $K_n$ that shows that
	$R(3,3,k) \geq n$. Let $G = E_R(\chi(K_n)) \cup E_B(\chi(K_n))$
	that is, $G$ is the graph of red and blue edges. Note that
	$\alpha(G) < k$, as the coloring had no green $K_k$.
	To apply [\ref{lemma:AKSz}], we must calculate $k_3(G)$ and we know
	it does not have any full red, or any full blue triangles.
	So let's bound the number of $\{R,R,B\}$ triangles in $G$ (the
	number of $\{B,B,R\}$ by simmetry will be the same).
	One way to count $\{R,R,B\}$ triangles is by picking the cherry vertex.
	\[
		k_{\{R,R,B\}}(G) = \sum_{v \in G} e_B(G[N_R(v)])
	\]
	By Erdös-Szekeres, we know that $N_R(v)$ is red-edge-free and
	$|N_R(v)| \leq R(3,k)$ for any vertex $v$.
	We also have that for any $u \in N_R(v)$, $N_B(u) \cap N_R(v)$ is a
	independent set, by the previous observation and the
	fact that we are blue triangle free. So $d_B[G[N_R(v)]](u) < k$
	(here this monstrosity means that the blue degree of $u$ in
	$N_R(v)$ is less than $k$).
	This implies that $e_B(G[N_R(v)]) < R(3,k)k/2$ for any $v$ and therefore that
	\[
		k_{\{R,R,B\}}(G) < \frac{nR(3,k)k}{2}
	\]
	so the total number of triangles is also of order $nk^3/\log k$.
	Applying [\ref{lemma:AKSz}] with $d = \frac{Ck^2}{\log k} \geq  R(3,k)$
	and this $k_3(G) < ndk$ we get that we can use $\lambda = Cd/k \geq
	\frac{Ck}{\log k}$. So we get
	\[
		k > \alpha(G) \geq \frac{n}{d} \log k \geq \frac{C' n (\log k)^2}{k^2}
	\]
	which implies $n \leq Ck^3/(\log k)^2$.
\end{proof}

Next class we will see a very cool lowerbound by Alon, who observed
the following lemma:
\begin{lemma}
	\label{lemma:few_indep_r(3,3,k)}
	If there exists $K_3$-free $G$ on $n$ vtx with number of
	independent $k$ sets of $G$ less than $\binom{n}{k}^{1/2}$
	then $R(3,3,k) > n$.
\end{lemma}
\begin{proof}
	Embed $G$ randomly in $K_n$ with color red, do it again in color
	blue. By the union bound, no independent $k$ set survive (so we win
	by painting non edges with green).
\end{proof}

We will turn for a pseudo-random construction to prove the existence
of a graph that satisfies [\ref{lemma:few_indep_r(3,3,k)}].
For this, we need a good definiton of what pseudo-random means.
\begin{definition}
	We say that a $n$-vtx graph $G$ is $(p,\beta)$-jumbled if $e(G) =
	p\binom{n}{2}$ and for every $S,T \subset V(G)$, we have that
	\[
		|e(S,T) - p|S||T|| \leq \beta\sqrt{|S||T|}.
	\]
	where here $e(S,T)$ is the number of edges that start in $S$ and end
	in $T$. Edges in $S \cap T$ are counted twice.
\end{definition}

The most pseudo-random graph (the best $\beta$ you can get) is with
$G(n,p)$, I don't know anyway to prove this now,
but let's calculate how jumbled $G(n,p)$ is.

\begin{theorem}
	$G(n,p)$ is $(p',\beta)$-jumbled with $p' = e(G(n,p))/\binom{n}{2} =
	(1 + o(1))p$ and $\beta = O(\sqrt{pn})$ w.h.p.
\end{theorem}
\begin{proof}
	We use Chernoff inequality [\ref{lemma:expec_binomial_bounds}], to
	note that for absolute $C,c > 0$ and taking $A$ suffciently big,
	\[
		\Pr( |e(S,T) - p|S||T|| \geq A\sqrt{pn|S||T|}) \leq C\exp(-cA^2n) \ll 4^{-n},
	\]
	so $G(n,p)$ is $\beta = A\sqrt{np}$-jumbled w.h.p by the union bound.
\end{proof}

This yields the following definition
\begin{definition}
	A graph is said to be \textbf{optimally pseudorandom} if it is
	$(p,\beta)$-jumbled with $\beta = O(\sqrt{pn})$.
\end{definition}

The following question is at the core of this method. How many
independent $k$-sets can a $(p,\beta)$-jumbled graph have? We now have
a deterministic fact about our graph and we can try to count. I won't
do it here, but the idea is to count sequences $(v_1, \dots, v_k)$
of vertices that can make a independent set.

Using the $(p,\beta)$
jumbled condition, we can count how many "bad" vertices there are
when trying to grow the sequence. We can make only a few
"good" steps (extension steps where our non-neighborhood decreases
with correct size),
because they reduce the non-neighborhood by a $(1-p)/2$ ratio which is just too
fast. And each "bad" step  (steps which don't reduce the non-neighborhood much)
we are obliged to choose from a small set. So counting
everything yields a result.

The following lemma will be stated without proof, it is an algebraic
construction.
\begin{lemma}
	(Quadrangle) There exists a $n$-vertex, $n^{1/3}$-regular,
	$n^{1/3}$-uniform hypergraph $H$ that satisfies:
	\begin{enumerate}[label=(\roman*)]
		\item $\forall e_1, e_2 \in E(H), \quad |e_1 \cap e_2| \neq 2$.
		\item $\forall e_1,e_2,e_3 \in E(H)$ it does not hold that $|e_1
			\cap e_2| = |e_2 \cap e_3| = |e_1 \cap e_3| = 1$.
	\end{enumerate}
\end{lemma}

Putting a randomly embedded $K_{t,t}$ with $t = n^{1/3}/2$ on each
edge we get a $K_3$-free, optimally pseudorandom graph
where $p$ is the maximum it can be, $p = n^{-1/3}$.
(As Rob pointed out an optimally pseudorandom, triangle-free graph
	must have density
at most $p = n^{-1/3}$).
